\documentclass[10pt]{article}
\usepackage{braket}
\usepackage{algorithm}
\usepackage{algpseudocode}
\usepackage{soul}
\usepackage{float}
\usepackage{tabularx}
\usepackage{mathtools}
\usepackage{amsmath, amssymb, amsthm}
\usepackage{tikz}
\usetikzlibrary{arrows}
\usetikzlibrary{decorations.text}
\usepackage{enumitem}
\usepackage{geometry}
\usepackage{fancyhdr}
\usepackage{titlesec}
\usepackage{xltabular}
\usepackage{color}
\usepackage{hyperref}
\geometry{margin=1in}

\pagestyle{fancy}
\fancyhf{}
\cfoot{\thepage}
\title{SFWRENG 3DX4 - Pre-lab 5}
\author{Matthew Farah, \href{mailto:farahm11@mcmaster.ca}{$\langle$farahm11@mcmaster.ca$\rangle$}, 400468251}
\date{\today}

\hypersetup{
	colorlinks=true,
	linkcolor=blue,
	filecolor=magenta,
	urlcolor=blue,
	citecolor=blue,
	anchorcolor=blue
}

\fancyhead{}
\fancyhead[L]{Matthew Farah, \href{mailto:farahm11@mcmaster.ca}{$\langle$farahm11@mcmaster.ca$\rangle$}, 400468251}
\fancyhead[R]{SFWRENG 3DX4 - Pre-lab 5}

\AtBeginEnvironment{align}{\setcounter{equation}{0}}

\newcommand{\comment}[1]{\parbox[t]{0.45\textwidth}{\raggedright #1}}

\setlength{\parindent}{0cm}

\setcounter{tocdepth}{3}
\hypersetup{bookmarksopen=true, bookmarksopenlevel=2, bookmarksdepth=3}

\newcounter{question}
\renewcommand{\thequestion}{\arabic{question}}

\newenvironment{question}{
	\refstepcounter{question}
	\phantomsection
	\addcontentsline{toc}{section}{\protect{\large{Question \thequestion}}}
	\vspace{1em}
	\par
	\large\textbf{Question \thequestion}
	\vspace{.2cm}
	\par
	\setcounter{subquestion}{0}
	\setcounter{step}{0}
}{\vspace{.5em}}

\newenvironment{solution}{
	\vspace{1em}\par\large\textbf{{Solution:}}\par\vspace{1em}
}{
	\vspace{1em}
	\newpage
}

\newcounter{subquestion}
\renewcommand{\thesubquestion}{\alph{subquestion}}

\newenvironment{subquestion}{
	\refstepcounter{subquestion}
	\phantomsection
	\vspace{1em}\par\large\textbf{(\thesubquestion)}\hspace{-.5em}
	\setcounter{step}{0}
	\addcontentsline{toc}{subsection}{\protect{\large{Part \thequestion.\thesubquestion}}}
}{
	\par
	\vspace{1em}
}

\makeatother

\makeatletter

\newcounter{step}
\renewcommand{\thestep}{\Roman{step}}

\newenvironment{step}{
	\refstepcounter{step}
	\vspace{1.5em}\par\noindent\large\textbf{(\thestep)}
}{
	\par
	\vspace{1.5em}
}

\makeatother

\newcolumntype{Y}{>{\centering\arraybackslash}X}
\renewcommand{\arraystretch}{1.8}

\fontsize{10pt}{14pt}\selectfont

\begin{document}

\maketitle

\tableofcontents
\newpage

\begin{question}
	In Lab 4, we used a PD compensator to control our ball and beam apparatus.
	The transfer function of our PD compensator was as follows:

	\begin{gather*}
		G_{c}(s) = \left( K_{D}s + K_{P} \right) \\
	\end{gather*}

	However, we did not use the compensator in this form. The transfer function
	we used in lab was as follows:

	\begin{gather*}
		G_{c}(s) = K_{C}\left(s + z \right) \\
	\end{gather*}

	Solve for $K_C$ and $z$ in terms of $K_D$ and $K_P$.
\end{question}

\begin{solution}

\end{solution}

\begin{question}
	Given that the transfer function of our Ball and Beam plant used in the previous
	lab is:

	\begin{gather*}
		G(s) = \frac{0.419}{s^2} \\
	\end{gather*}

	And given that the controller is applied to the plant in cascade configuration,
	find: \\
\end{question}

\begin{subquestion}
	Static error constant for position (position constant).
\end{subquestion}

\begin{solution}

\end{solution}

\begin{subquestion}
	Static error constant for velocity (velocity constant).
\end{subquestion}

\begin{solution}

\end{solution}

\begin{subquestion}
	Static error constant for acceleration (acceleration constant).
\end{subquestion}

\begin{solution}

\end{solution}

\begin{subquestion}
	Static error constant for step input $u(t)$.
\end{subquestion}

\begin{solution}

\end{solution}

\begin{subquestion}
	Static error constant for ramp input $tu(t)$.
\end{subquestion}

\begin{solution}

\end{solution}

\begin{subquestion}
	Static error constant for parabolic input $t^2u(t)$.
\end{subquestion}

\begin{solution}

\end{solution}

\begin{question}
	In Lab 5, we will be augmenting our controller to include an integrator. The
	transfer function of our new PID compensator will be as follows:

	\begin{gather*}
		G_{c}(s) = \left(K_{D}s + K_{P} + K_{I}s \right) \\
	\end{gather*}

	Given that the transfer function for our plant has not changed, and given that
	this controller is also applied to the plant in cascade configuration, find:
\end{question}

\begin{subquestion}
	Static error constant for position (position constant).
\end{subquestion}

\begin{solution}

\end{solution}

\begin{subquestion}
	Static error constant for velocity (velocity constant).
\end{subquestion}

\begin{solution}

\end{solution}

\begin{subquestion}
	Static error constant for acceleration (acceleration constant).
\end{subquestion}

\begin{solution}

\end{solution}

\begin{subquestion}
	Static error constant for step input $u(t)$.
\end{subquestion}

\begin{solution}

\end{solution}

\begin{subquestion}
	Static error constant for ramp input $tu(t)$.
\end{subquestion}

\begin{solution}

\end{solution}

\begin{subquestion}
	Static error constant for parabolic input $t^2u(t)$.
\end{subquestion}

\begin{solution}

\end{solution}

\begin{question}
	Consider the block diagram shown below where $D(s)$ is a step disturbance input. \\

	%TODO: Add block diagram.

	Ideally you want your controller design to reject a step disturbance input at
	$D(s)$. This means that in the steady state for $D(s) = \frac{1}{s}$, the value of $Y(s)$ is
	unchanged.
\end{question}

































\end{document}
